\chapter*{Abstract}
\addcontentsline{toc}{chapter}{Abstract}
The ADNEX model, introduced in 2014, has demonstrated good diagnostic ability for five ovarian tumor types: benign, borderline, stage I invasive, stage II-IV invasive, and secondary metastasis. However, ADNEX should be further improved based on previous findings. First, despite excellent overall performance, we should understand and reduce observed differences between centers. Second, ADNEX was developed on patients selected for surgery, and should be adapted to work for all patients with a newly detected tumor. Third, we need to understand whether ADNEX works less well when used by clinicians with limited experience. Fourth, we should compare different statistical algorithms for developing the model. Finally, ADNEX is the fruitful result of a combination of clinical and statistical research. Therefore, we need to continue the statistical research on prediction modeling. In particular, we are interested in development and validation of machine learning models in the presence of multicenter (i.e. clustered) data, and in variable selection methods based criteria for clinical utility rather than for statistical performance.